\documentclass{amsart}
% For dealing with references we use the comment environment
\usepackage{verbatim}
\newenvironment{reference}{\comment}{\endcomment}
%\newenvironment{reference}{}{}
\newenvironment{slogan}{\comment}{\endcomment}
\newenvironment{history}{\comment}{\endcomment}

% For commutative diagrams we use Xy-pic
\usepackage[all]{xy}

% We use 2cell for 2-commutative diagrams.
\xyoption{2cell}
\UseAllTwocells

% We use multicol for the list of chapters between chapters
\usepackage{multicol}

% This is generally recommended for better output
\usepackage{lmodern}
\usepackage[T1]{fontenc}

% For cross-file-references

% Package for hypertext links:
\usepackage{hyperref}

% For any local file, say "hello.tex" you want to link to please

% Theorem environments.
%
\theoremstyle{plain}
\newtheorem{theorem}[subsection]{Theorem}
\newtheorem{proposition}[subsection]{Proposition}
\newtheorem{lemma}[subsection]{Lemma}

\theoremstyle{definition}
\newtheorem{definition}[subsection]{Definition}
\newtheorem{example}[subsection]{Example}
\newtheorem{exercise}[subsection]{Exercise}
\newtheorem{situation}[subsection]{Situation}

\theoremstyle{remark}
\newtheorem{remark}[subsection]{Remark}
\newtheorem{remarks}[subsection]{Remarks}

\numberwithin{equation}{subsection}

% Macros
%
\def\lim{\mathop{\mathrm{lim}}\nolimits}
\def\colim{\mathop{\mathrm{colim}}\nolimits}
\def\Spec{\mathop{\mathrm{Spec}}}
\def\Hom{\mathop{\mathrm{Hom}}\nolimits}
\def\Ext{\mathop{\mathrm{Ext}}\nolimits}
\def\SheafHom{\mathop{\mathcal{H}\!\mathit{om}}\nolimits}
\def\SheafExt{\mathop{\mathcal{E}\!\mathit{xt}}\nolimits}
\def\Sch{\mathit{Sch}}
\def\Mor{\mathop{\mathrm{Mor}}\nolimits}
\def\Ob{\mathop{\mathrm{Ob}}\nolimits}
\def\Sh{\mathop{\mathit{Sh}}\nolimits}
\def\NL{\mathop{N\!L}\nolimits}
\def\CH{\mathop{\mathrm{CH}}\nolimits}
\def\proetale{{pro\text{-}\acute{e}tale}}
\def\etale{{\acute{e}tale}}
\def\QCoh{\mathit{QCoh}}
\def\Ker{\mathop{\mathrm{Ker}}}
\def\Im{\mathop{\mathrm{Im}}}
\def\Coker{\mathop{\mathrm{Coker}}}
\def\Coim{\mathop{\mathrm{Coim}}}

% Boxtimes
%
\DeclareMathSymbol{\boxtimes}{\mathbin}{AMSa}{"02}

%
% Macros for moduli stacks/spaces
%
\def\QCohstack{\mathcal{QC}\!\mathit{oh}}
\def\Cohstack{\mathcal{C}\!\mathit{oh}}
\def\Spacesstack{\mathcal{S}\!\mathit{paces}}
\def\Quotfunctor{\mathrm{Quot}}
\def\Hilbfunctor{\mathrm{Hilb}}
\def\Curvesstack{\mathcal{C}\!\mathit{urves}}
\def\Polarizedstack{\mathcal{P}\!\mathit{olarized}}
\def\Complexesstack{\mathcal{C}\!\mathit{omplexes}}
% \Pic is the operator that assigns to X its picard group, usage \Pic(X)
% \Picardstack_{X/B} denotes the Picard stack of X over B
% \Picardfunctor_{X/B} denotes the Picard functor of X over B
\def\Pic{\mathop{\mathrm{Pic}}\nolimits}
\def\Picardstack{\mathcal{P}\!\mathit{ic}}
\def\Picardfunctor{\mathrm{Pic}}
\def\Deformationcategory{\mathcal{D}\!\mathit{ef}}

\setlength{\emergencystretch}{3em}
\begin{document}
\title[Stacks Project, Tag 0CD4]{570 --- Smooth proper curves of genus at least two have separable closed points of degree at most twice the genus minus two}
\maketitle
\noindent Copyright (C) 2005--2025 Johan de Jong. Stacks Project authors. Source: \href{https://stacks.math.columbia.edu/tag/0CD4}{Tag 0CD4}.

\noindent Editorial difficulty note (not author text). Difficulty H1: The usual canonical pencil yields a degree bound but may be inseparable over an imperfect ground field. The proof rules out its Frobenius-factorization alternative using linear-series bounds on the Frobenius twist, then chooses a rational point in the étale locus of the resulting pencil..

\noindent Editorial provenance note (not author text). History: excerpt from revision \texttt{a0444\allowbreak 6e57e\allowbreak c1fbc\allowbreak 252a8\allowbreak 71afc\allowbreak ec775\allowbreak 2fb28\allowbreak 07b14}, curves.tex, lines 2753--2811. Reference presentation was adapted; original-author context and core support blocks were retained or added. The mathematical wording is unchanged. Licensed under GNU FDL 1.2 or later; no invariant sections or cover texts. See ../licenses/COPYING. Context blocks reproduce author definitions and supporting results; they are not additional counted problems.

\begin{lemma}
\label{lemma-point-over-separable-extension}
Let $k$ be a field. Let $X$ be a smooth proper curve over $k$
with $H^0(X, \mathcal{O}_X) = k$ and genus $g \geq 2$.
Then there exists a closed point $x \in X$ with
$\kappa(x)/k$ separable of degree $\leq 2g - 2$.
\end{lemma}

\begin{proof}
Set $\omega = \Omega_{X/k}$. By
Lemma \href{https://stacks.math.columbia.edu/tag/0C1A}{lemma genus smooth (Tag 0C1A)} this has degree $2g - 2$
and has $g$ global sections. Thus we have a $\mathfrak g^{g - 1}_{2g - 2}$.
By the trivial Lemma \href{https://stacks.math.columbia.edu/tag/0CCQ}{lemma linear series trivial existence (Tag 0CCQ)}
there exists a $\mathfrak g^1_{2g - 2}$
and by Lemma \href{https://stacks.math.columbia.edu/tag/0CCR}{lemma g1d (Tag 0CCR)} we obtain a morphism
$$
\varphi : X \longrightarrow \mathbf{P}^1_k
$$
of some degree $d \leq 2g - 2$. Since $\varphi$ is flat
(Lemma \href{https://stacks.math.columbia.edu/tag/0CCK}{lemma flat (Tag 0CCK)}) and finite
(Lemma \href{https://stacks.math.columbia.edu/tag/0CCL}{lemma finite (Tag 0CCL)})
it is finite locally free of degree $d$
(Morphisms, Lemma \href{https://stacks.math.columbia.edu/tag/02KB}{lemma finite flat (Tag 02KB)}).
Pick any rational point $t \in \mathbf{P}^1_k$
and any point $x \in X$ with $\varphi(x) = t$.
Then
$$
d \geq [\kappa(x) : \kappa(t)] = [\kappa(x) : k]
$$
for example by
Morphisms, Lemmas \href{https://stacks.math.columbia.edu/tag/0CC2}{lemma finite locally free universally bounded (Tag 0CC2)}
and \href{https://stacks.math.columbia.edu/tag/03J5}{lemma characterize universally bounded (Tag 03J5)}.
Thus if $k$ is perfect (for example has characteristic zero
or is finite) then the lemma is proved. Thus we reduce to the
case discussed in the next paragraph.

\medskip\noindent
Assume that $k$ is an infinite field of characteristic $p > 0$.
As above we will use that $X$ has a $\mathfrak g^{g - 1}_{2g - 2}$.
The smooth proper curve $X^{(p)}$ has the same genus as $X$.
Hence its genus is $> 0$. We conclude that $X^{(p)}$ does not have a
$\mathfrak g^{g - 1}_d$ for any $d \leq g - 1$ by
Lemma \href{https://stacks.math.columbia.edu/tag/0CCS}{lemma grd inequalities (Tag 0CCS)}.
Applying Lemma \href{https://stacks.math.columbia.edu/tag/0CD3}{lemma inseparable linear system (Tag 0CD3)}
to our $\mathfrak g^{g - 1}_{2g - 2}$ (and noting that $2g - 2/p \leq g - 1$)
we conclude that possibility (2) does not occur. Hence we obtain a morphism
$$
\varphi : X \longrightarrow \mathbf{P}^1_k
$$
which is generically \'etale (in the sense of the lemma)
and has degree $\leq 2g - 2$. Let $U \subset X$ be the nonempty
open subscheme where $\varphi$ is \'etale. Then
$\varphi(U) \subset \mathbf{P}^1_k$ is a nonempty Zariski open
and we can pick a $k$-rational point $t \in \varphi(U)$ as $k$ is infinite.
Let $u \in U$ be a point with $\varphi(u) = t$.
Then $\kappa(u)/\kappa(t)$ is separable
(Morphisms, Lemma \href{https://stacks.math.columbia.edu/tag/02GL}{lemma etale over field (Tag 02GL)}),
$\kappa(t) = k$, and $[\kappa(u) : k] \leq 2g - 2$ as before.
\end{proof}
\end{document}
